Beide Seiten einer Gleichung haben dieselbe Einheit, ebenso alle Summanden.
\begin{abcliste}
\abc $p = \dfrac{v}{s}$ hat die Einheit $\dfrac{\si{\meter\per\second}}{\si{\meter}}$, also $[p] = \result{\si{\per\second}}$.
\abc $q = \dfrac{v\,t^2}{s}$ hat die Einheit $\dfrac{\si{\meter\per\second}\cdot\si{\second\squared}}{\si{\meter}}$, also $[q] = \result{\si{\second}}$.
\abc $p\cdot v$ muss die Einheit Meter haben: $[p] = \dfrac{\si{\meter}}{\si{\meter\per\second}} = \result{\si{\second}}$. Der Summand $q$ ist eine Strecke: $[q] = \result{\si{\meter}}$.
\abc $\left(p\,v\right)^2 = \dfrac{s}{t}$ hat die Einheit \si{\meter\per\second}, also hat $p\,v$ die Einheit $\left(\si{\meter\per\second}\right)^{1/2}$ und
\begin{align}
 [p] &= \frac{\left(\si{\meter\per\second}\right)^{1/2}}{\si{\meter\per\second}} = \result{\left(\si{\second\per\meter}\right)^{1/2}}
\end{align}
\abc $\dfrac{p}{v}$ ist eine Strecke: $[p] = \si{\meter}\cdot\si{\meter\per\second} = \result{\si{\meter\squared\per\second}}$. Ebenso $q\,t^2$: $[q] = \result{\si{\meter\per\second\squared}}$.
\end{abcliste}
